\section{A. The untwisted limit}
\subsection{A.1. The nodal plumbing}

We write $L=\mathcal O_S(P-O)$, $T=t_{6P}$,
$J=L\otimes T^*L^{-1}=\pi^*\mathcal O_B(1)$, and
$R=\mathbb P_S(\mathcal O\oplus L)$. The convention for this projective
bundle is that its zero and infinity sections have normal bundles
$L$ and $L^{-1}$. The identification defining $Q$ is
$s_0(x)\sim s_\infty(Tx)$. The equality defining $J$ is the equality
in Engel's Lemma 2.1, with this direction of $T$; its intersection
calculation uses $P\cdot O=0$, $P^2=O^2=-1$,
$6P\cdot O=6P\cdot P=2$. Thus the degree on $O$ is
$(O-P)\cdot6P+(P-O)\cdot O=1$. We use the section
$\phi\in H^0(S,J)$ whose divisor is $S_\infty$, as in his
Corollary 2.2. The source being checked is the attached version of
\cite{E26}, not a later inference from its conjecture.

\begin{proposition}[A.1.1; Proved]\label{push:A11}
The cones $P_{k,l}$ in Lead 2 and their faces form a fan. Its two
period automorphisms on the cocharacter lattice are
\[
 G_1(y_1,y_2,h_q,h_u)=(y_1+h_q,y_2,h_q,h_u),\qquad
 G_2(y_1,y_2,h_q,h_u)=(y_1,y_2+h_q+h_u,h_q,h_u).
\]
The affine chart of $P_{0,0}$ is
\[
 U_{0,0}=\operatorname{Spec}\mathbb C[x,y,z,v,u]/(xy-uzv),
 \qquad q=zv.
\]
In particular the unsubdivided chart is not a smooth simultaneous
model of the cusp.
\end{proposition}
\begin{proof}
At fixed $h_q,h_u\geq0$, the cone $P_{k,l}$ cuts out the rectangle
\[
 kh_q\leq y_1\leq(k+1)h_q,\qquad
 l(h_q+h_u)\leq y_2\leq(l+1)(h_q+h_u).
\]
When $h_q=0$ the first interval is a point. These rectangles tile
their support, and their intersections are common faces, including
this collapsed boundary. This proves the fan assertion. The displayed
integral transformations preserve the fan and have integral inverses.

For a character $(a,b,c,d)$, the six inequalities of the dual cone
are $c,a+c,b+c,a+b+c,d,b+d\geq0$. The characters
\[
 x=(0,1,0,0),\quad y=(0,-1,1,1),\quad
 z=(1,0,0,0),\quad v=(-1,0,1,0),\quad u=(0,0,0,1)
\]
generate its integral monoid. Indeed if $b\geq0$ we subtract $b x$;
if $b<0$ we subtract $(-b)y$. The remaining character has second
coordinate zero and is a nonnegative integral combination of $z,v,u$.
The sole relation is $x+y=u+z+v$, and the character of the first
base coordinate is $z+v$. This proves the coordinate-ring assertion,
including normality. The calculation is checked exactly in
\texttt{code/ab\_push\_local\_checks.py}.
\end{proof}

\begin{theorem}[A.1.2; Proved]\label{push:A12}
There is a reduced complex space $\mathcal X^\flat$ and a proper flat
holomorphic map
\[
 (f,\varpi):\mathcal X^\flat\longrightarrow B\times\Delta_u
\]
whose central threefold is $Q$. For $u\ne0$ its fibre is the
contraction of Engel's chosen $X(u)$ at one rational curve of the
Mumford fibre. This curve is the double curve corresponding to the
diagonal of the square in the chosen $A_2$ triangulation; in the
return's coordinates its direction is $e_1-e_2$. Its normal bundle
in $X(u)$ is $\mathcal O_{\mathbb P^1}(-1)^{\oplus2}$. Its image is
the unique ordinary double point of $\mathcal X^\flat_u$.
\end{theorem}
\begin{proof}
Here is a construction on the whole compact surface $S$. Choose a
Hermitian metric on $L$, use the dual metric on $L^{-1}$, and choose
$0<r<1/4$. For sufficiently small $\epsilon>0$ we have
$|u|\|\phi_x\|<r^2$ for all $x\in S$ and $|u|<\epsilon$.
Take the core
\[
 C=\{(x,a,u):a\in L_x,\ r<\|a\|<r^{-1},\ |u|<\epsilon\}
\]
and the node space
\[
 N=\{(x,a,b,u):a\in L_x,\ b\in L^{-1}_{Tx},\
 a\otimes b=u\phi_x,\ \|a\|,\|b\|<2r,\ |u|<\epsilon\}.
\]
On $r<\|a\|<2r$ identify a point of $N$ with $(x,a,u)\in C$.
On $r<\|b\|<2r$ identify it with $(Tx,b^{-1},u)\in C$.
Both identifications are biholomorphisms between open complex spaces;
in the first collar $b=u\phi_x/a$, and in the second
$a=u\phi_xb^{-1}$. Their domains are disjoint because
$\|a\|\|b\|<r^2$. Changes of frames multiply the equation by
units and give the stated bundle-valued equation, so no choice of
frame enters the gluing.

This gluing is Hausdorff. At either end of a collar, exactly one
of its two charts contains the limiting circle: the core omits
$\|a\|=r$, and the node chart omits $\|a\|=2r$; the infinity
collar has the identical property in $b$. Thus the graph of the
identification has no unpaired pair of limiting points. The two
collars in each chart are separated, so there are no additional
chains of identifications. These observations give disjoint
saturated neighbourhoods of any two inequivalent points. Properness
is equally concrete. Over a compact subset of $B\times\Delta_\epsilon$,
the closed core $3r/2\leq\|a\|\leq2/(3r)$ together with the
closed node subspace $\|a\|,\|b\|\leq7r/4$ is compact and maps
onto the glued space. To see surjectivity, a point outside the smaller
core lies on a node collar, and a node point outside the smaller
node subspace lies on the smaller core. Compactness of $S$ supplies
the uniform bound used above. This proves properness of the joint
map to $(t,u)$.

The total node rings are hypersurface rings over the regular
two-dimensional local rings of $S$. They are Cohen--Macaulay.
Every fibre of the joint map has pure dimension two: first fix
$t$, and then impose $xy=u\phi$ on the two line coordinates over
the one-dimensional fibre of $S$. Consequently $t-t_0,u-u_0$ is a
regular sequence in every total local ring. The regular-sequence
flatness criterion over the smooth two-dimensional base proves
flatness. On the core the same conclusion is immediate from
flatness of $S\to B$. This includes the nonreduced Kodaira fibres
at $0,1$; smoothness of $S$, rather than smoothness of those fibres,
is what the argument uses. At $u=0$ the node equation is $ab=0$,
and the two collar maps recover precisely the specified pinching
of the two sections of $R$.

For $u\ne0$ and $t\ne\infty$, the equation gives
$a_{Tx}=a_x/(u\phi_x)$. Hence the resulting quotient is exactly
$L^\times/\langle u\phi\rangle$, where the contracting generator
covers $T^{-1}$. Every orbit meets a core or a collar, by the
uniform contraction estimate; the displayed comparison identifies
exactly its orbit relation. It is therefore a biholomorphism,
not just a meromorphic comparison.

At the nodal point of $S_\infty$, write $q=zv$. Absorbing a unit
in one line coordinate gives $xy=uzv$. At a smooth point of
$S_\infty$ the equation is $xy=uq$. For $u\ne0$ the latter
threefold is smooth; the former has one nondegenerate quadratic
singularity at the origin. There are no other singularities of a
nonzero $u$ fibre.

For completeness, the relation to the infinite fan quotient is
also explicit. In a Tate covering of a neighbourhood of the $I_1$
fibre use the torus coordinates $(Z,W)$ of Engel's Proposition 2.7.
The two periods act as
\[
 (Z,W,q,u)\longmapsto(c_1qZ,c_2W,q,u),\qquad
 (Z,W,q,u)\longmapsto(c_3Z,c'_4uqW,q,u),
\]
with holomorphic units $c_i$. Their monomial parts are $G_1,G_2$.
The chart indexed by $(k,l)$ has
$z=Zq^{-k}$, $v=q^{k+1}/Z$, $x=W(qu)^{-l}$,
$y=qu/x$, with the harmless unit changes supplied by the line
frames. It is the chart of Proposition~\ref{push:A11}.
The Tate cover unfolds the node of $S_\infty$ into its infinite
chain, and unfolding the two plumbing collars gives the additional
integer index $l$. The foregoing collar identifications are exactly
the two deck identifications. Thus, on an invariant neighbourhood
over a sufficiently small bidisc, the fan space maps by a
holomorphic covering to the plumbing, with deck group $\mathbb Z^2$.
Each core or node coordinate neighbourhood has disjoint lifts
indexed by the two unfolded integers; this proves freeness and
proper discontinuity, including the $\rho_l$ orbits. The first
period stabilizes such a ray but on its orbit translates the Tate
coordinate by $c_1q$, with $q\ne0$ and $|c_1q|<1$; it does not
fix a point. Properness of its quotient was proved by the compact
core and node representatives. This supplies the analytic
comparison needed here. Construction 3.8 of \cite{EGFS25} is a
related multivariable polarized construction; its symmetry
hypothesis on period units is not being silently imposed on these
generally unpolarized tori.

Subdividing the square at $h_u=0$ by the chosen diagonal gives
the usual small resolution of $xy=uzv$ for $u\ne0$. Its exceptional
curve has the two affine charts of the total space of
$\mathcal O_{\mathbb P^1}(-1)\oplus
\mathcal O_{\mathbb P^1}(-1)$, as follows by setting $x=za$ on
one chart and $z=xb$ on the other. The two normal coordinates
transform by the reciprocal affine coordinate. There is one
diagonal edge modulo the periods, so there is exactly one such
curve. Removing that edge contracts the paired opposite
$(-1)$-curves in the normalization $dP_6$ to the square model
$\mathbb P^1\times\mathbb P^1$. This identifies the contracted
curve intrinsically in the selected toric model. It is not the
$e_2$-direction curve fixed by $\operatorname{Aut}^0$.
\end{proof}

\subsection{A.2. Two minimal toroidal fillings and further modifications}

\begin{proposition}[A.2.1; Proved]\label{push:A21}
Among subdivisions with no new rays, the fan in A.1.1 has exactly
two smooth period-invariant subdivisions restricting to the chosen
$A_2$ triangulation at $h_u=0$. These are also all subdivisions
which are flat over the original $(q,u)$ bidisc, have reduced
$u$-central divisor, and retain that boundary triangulation.
Without the joint-flatness restriction there are infinitely many
smooth invariant subdivisions retaining that triangulation and
having reduced $u$-central divisor.
\end{proposition}
\begin{proof}
Write the six rays of $P_{0,0}$ as
$A=(0,0,1,0)$, $B=(1,0,1,0)$, $C=(0,1,1,0)$,
$D=(1,1,1,0)$, $E=(0,0,0,1)$, $F=(0,1,0,1)$.
The return's $A_2$ diagonal is $BC$. A triangular prism has six
triangulations without new vertices. With this bottom diagonal
the three possibilities are
\[
\begin{array}{c|c|c}
\text{maximal tetrahedra}&\text{left diagonal}&\text{right diagonal}\\\hline
ABCE,\ BCDE,\ CDEF&CE&DE\\
ABCE,\ BCDF,\ BCEF&CE&BF\\
ABCF,\ ABEF,\ BCDF&AF&BF.
\end{array}
\]
Periodicity identifies the left and right faces by $A\mapsto B$,
$C\mapsto D$, $E\mapsto E$, $F\mapsto F$; it excludes exactly
the middle row. Every listed maximal cone has determinant $\pm1$.
There are no missed triangulations: the prism has normalized volume
three, all its nondegenerate four-vertex tetrahedra have normalized
volume one, and enumeration of their triples with pairwise
face intersections gives exactly these six possibilities before
fixing the bottom. The program \texttt{ab\_push\_fan.py} performs
this enumeration in rational arithmetic; intersections are tested
by all vertices of the combined barycentric halfspaces.

For the stronger assertion, a new ray with $h_q,h_u>0$ would
give a divisor mapping to the origin of the two-dimensional base.
Its fibre dimension would be three, contradicting joint flatness.
A ray with $h_u=0$ is excluded by the retained boundary
triangulation. On $h_q=0$, the fan support has $y_1=0$.
Reducedness requires the primitive ray to have $h_u=1$, so it
is already one of the $\rho_l$. Thus there are no new rays under
the stated three restrictions. Conversely the two triangulations
have all heights zero or one, and every cone maps onto a face of
the quadrant with saturated image. Their local monomial maps
are flat and their $u$ divisors reduced.

If joint flatness is omitted, star-subdivide the orbit of a
regular two-face generated by a bottom ray and a $\rho$ ray,
using their sum. The new ray has $h_u=1$; regularity is preserved
because replacing either basis vector by their sum has determinant
one. It leaves the bottom unchanged. Repeating the operation on
the resulting mixed faces gives infinitely many invariant smooth
refinements, all reduced over $u=0$, with increasing numbers of
rays modulo the period group. They are not jointly flat. This
distinction is necessary in reading ``all subdivisions'' in A.2.
\end{proof}

\begin{theorem}[A.2.2; Proved]\label{push:A22}
The two minimal subdivisions of A.2.1 give proper flat holomorphic
families $\mathcal X^+,\mathcal X^-\to\Delta_u$ with smooth total
spaces. Their joint maps to $B\times\Delta_u$ are flat. Over
$\Delta_u^*$ they are isomorphic as families to the specified
untwisted family $X(u)$. Each central fibre $Q^\pm$ is reduced,
irreducible, and has normal crossings, with smooth double locus
$S$ and no triple points.

Its normalization $R^\pm$ is obtained by blowing up $R$ along
$S_\infty$ in, respectively, one of its two sections and then
taking the small resolution at the unique ordinary double point
of that blow-up prescribed by the fixed bottom diagonal.
The two proper transforms of the sections are still copies of
$S$, glued by $T$. Their normal-bundle product is trivial;
in particular $\mathcal T^1_{Q^\pm}\simeq\mathcal O_S$.
\end{theorem}
\begin{proof}
Apply either fan subdivision over the cusp neighbourhood in
Theorem~\ref{push:A12}, and use the identity modification elsewhere.
Toric modifications are proper, and a subdivision is an
isomorphism off the strata it subdivides. Its restriction over
$q\ne0$ is the identity, so this gluing has no additional
transition choice. The period actions preserve the subdivisions,
and the covering description in A.1.2 descends the proper
modifications to the quotient. All maximal cones are regular.
Every local ring of the resulting fourfold is therefore regular.
The coordinate $u$ is the product of those affine coordinates
whose rays have $h_u=1$; there are at most two such rays in a
maximal cone. This proves reduced normal crossings and the absence
of triple points for the $u$-central fibre. The regular-sequence
argument, or the saturated face images in A.2.1, proves joint
flatness. Properness follows from properness of the toric
modification and of $\mathcal X^\flat$.

We give the normalization calculation, including the singular
centre. At a node of $S_\infty$, write the original normalization
as coordinates $(X,z,v)$, with section $X=0$ and centre
$(X,zv)$. Blowing up that ideal has charts
\[
 X=zvA;\qquad zv=XB.
\]
The first is smooth, and the second has one conifold at its
origin. A small resolution of the latter is, for one of the
bottom-diagonal conventions,
\[
 X=za,\quad v=aB;\qquad z=Xd,\quad B=dv.
\]
Interchanging $z,v$ gives the other diagonal convention. These
are holomorphic charts of a proper small modification: they are
the two usual charts of the blow-up of the appropriate rank-one
ideal on the conifold. The strict transform of $X=0$ is $A=0$
in the first chart and never passes through the conifold. It is
therefore unchanged as a surface. The other section is disjoint
from the centre and is unchanged as well. Over a smooth point
of $S_\infty$ the exceptional fibre is one $\mathbb P^1$;
over its node it is a chain of two $\mathbb P^1$'s, one being
the exceptional curve of the small resolution. This identifies
the normalization morphism, including its order of operations.

For example the corresponding total fourfold modification of
$xy=uzv$ can be written by first taking $x=ua$, $ay=zv$ and
then resolving the conifold in $(a,y,z,v)$, together with the
other chart $u=xb$, $y=bzv$. It has three smooth affine charts.
The assignments in \texttt{ab\_push\_local\_checks.py} verify
the equations and the map to $q=zv$ exactly. Exchanging the
chosen zero and infinity branches gives the other model.

The normal bundle of the section containing the centre changes
by tensoring with $\mathcal O_S(-S_\infty)$; the second normal
bundle is unchanged. Since the original product is $J$, the
new product is $J(-S_\infty)=\mathcal O_S$. The local
formula for $\mathcal T^1$ is the product of the two normal
bundles, with $T$ used on the infinity branch. This proves
d-semistability in this precise sense. The $\rho_l$ form one
orbit, so the central fibre is irreducible and its conductor is
exactly the two stated sections. Over $u\ne0$ both subdivisions
restrict to the fixed bottom triangulation, giving precisely
Engel's Mumford filling, simultaneously in $u$.

The birational comparison $\mathcal X^+\dashrightarrow
\mathcal X^-$ is the identity over the punctured $u$ disc and
off the exceptional loci over $(q,u)=(0,0)$. Locally it replaces
the two compatible side-face diagonals by the other two.
All its rays lie in the Gorenstein hyperplane $h_q+h_u=1$,
so both small modifications are crepant over
$\mathcal X^\flat$. Over a smooth point of the nodal curve
this is the ordinary Atiyah flop times a curve; at the node its
extension is exactly the displayed pair of prism subdivisions.
This specifies the birational map without asserting a uniqueness
theorem for arbitrary non-toroidal fillings.
\end{proof}

\begin{problem}[A.2.3; Open; unrestricted scope]\label{push:A23}
Classify all period-invariant smooth
subdivisions retaining the given bottom
triangulation, without imposing joint
flatness, and describe every resulting
central normalization and gluing.
A.2.1--A.2.2 classify and construct the
two minimal, jointly flat, reduced
models. They do not enumerate arbitrary
further subdivisions. Even within the
reduced class there are infinitely many,
as A.2.1 proves. If reducedness is also
omitted, new primitive rays may have
$h_u>1$, producing central multiplicities.
The literal unrestricted scope of
prompt A.2 is therefore not closed by
the two-model theorem.
\end{problem}

\subsection{A.3. The virtual Euler obstruction and the conjecture}

\begin{lemma}[A.3.1; Proved]\label{push:A31}
Let $Z$ be a compact reduced local-complete-intersection complex
space of pure dimension $n$. In any proper flat holomorphic
family over a disc with central fibre $Z$ and smooth nearby
fibres, their Euler number is
\[
 e_{\mathrm{vir}}(Z)=\langle c_n(T_Z^{\mathrm{vir}}),[Z]\rangle.
\]
This statement needs no projectivity or K\"ahler hypothesis.
\end{lemma}
\begin{proof}
Flat deformation preserves the local-complete-intersection
property relative to the disc near $Z$. Locally this follows by
lifting an embedding and its regular sequence of equations;
Nakayama's lemma and flatness show that their lifts generate
the ideal and remain a regular sequence. Properness allows one
to shrink the disc so that these neighbourhoods cover the total
space. The relative cotangent complex is consequently perfect
of amplitude $[-1,0]$. Its dual defines a topological complex
$K$-class $V$ on that total space; local finite vector-bundle
complexes define the class by descent, or by stabilization and
partitions of unity; the analytic-to-topological comparison is
given in \cite[Lemmas 3.1--3.2 and Section 3.1]{Levy08}.
Its restriction to a smooth fibre is its
tangent bundle, and its restriction to $Z$ is its intrinsic
virtual tangent class.

The compact fundamental cycles of the central and a nearby
fibre are homologous in the total space. More explicitly,
the inverse image $K$ of a short closed radial segment
$[0,\epsilon]$ is compact and semianalytic. Away from
$\operatorname{Sing}Z$ it is an oriented real
$(2n+1)$-manifold with boundary the nearby fibre and
$-Z_{\mathrm{reg}}$, because flatness with smooth fibre
implies smoothness of the morphism at those points.
The exceptional set $\operatorname{Sing}Z$ has real
dimension at most $2n-2$. Triangulate compatibly with
this set and the two endpoints, using
\cite[Section 3, Theorems 1--2]{Loj64}. The sum of the
oriented top simplices has boundary exactly
$[X_\epsilon]-[Z]$: no real $2n$-simplex can be supported
on that exceptional set. Reducedness and flatness give
multiplicity one on every central component. Pairing the globally defined
$c_n(V)$ with these homologous cycles gives the equality of
Chern numbers. Gauss--Bonnet on the smooth compact complex
fibre gives its Euler number. This proves the analytic assertion
directly. The algebraic counterpart is Theorem 5.8 and
Remark 5.9 of \cite{FG10}; those scheme statements alone would
not justify omitting the analytic comparison above.
\end{proof}

\begin{theorem}[A.3.2; Proved; Conjecture 8.2 refuted]\label{push:A32}
For the specified pinching $Q$ one has $e_{\mathrm{vir}}(Q)=0$.
Every smooth proper-flat smoothing of $Q$ therefore has Euler
number zero. In particular there is no proper flat holomorphic
filling of Engel's untwisted family with central fibre exactly
$Q$, even if no map to $B$ is required. Thus Conjecture 8.2
of the attached \cite{E26} is false as worded. The proper nodal
filling of A.1.2 and the two natural semistable fillings of
A.2.2 are precise replacements, with different central fibres.
\end{theorem}
\begin{proof}
Let $D_0,D_\infty\subset R$ be the conductor sections and
$\nu:R\to Q$ the normalization. On either branch of the
hypersurface $xy=0$, the derivative comparison with its
normalization gives in $K^0(R)$, where pullback of the
perfect tangent class means derived pullback,
\[
 \nu^*T_Q^{\mathrm{vir}}
 =T_R-i_{0*}(J|_{D_0})-i_{\infty*}(J|_{D_\infty}).
\]
Indeed in a smooth local ambient space the normal line $A$ of
the branch and the normal line of the whole hypersurface
differ by $D$: the latter is $A(D)$. Their difference in
$K$-theory is $-[i_*(A(D)|_D)]$, and $A(D)|_D$ is exactly
the product of the two branch normals, namely $\mathcal T^1_Q$.
These local comparisons are intrinsic and agree on overlaps.
On the infinity branch $J$ means its $T^{-1}$-pullback
from the double locus, canonically the same line bundle
because it comes from $B$. Thus $J$ extends from both sections by the pullback from $B$,
so $[i_*J]=[J]-[J(-D)]$.

Put $j=c_1(J)=\pi_R^*F$, $d_i=[D_i]$. Since the sections are
disjoint, $d_0d_\infty=0$, and since $F$ is a fibre, $j^2=0$.
Consequently
\[
 c(\nu^*T_Q^{\mathrm{vir}})
 =c(T_R)\frac{(1+j-d_0)(1+j-d_\infty)}{(1+j)^2}
 =c(T_R)\{1-d_0-d_\infty+j(d_0+d_\infty)\}.
\]
The degree-three part has integral
\[
 e(R)-\sum_i\int_{D_i}c_2(T_R|_{D_i})
       +\sum_i\int_{D_i}j\,c_1(T_R|_{D_i}).
\]
We have $e(R)=2e(S)=24$. On a section with normal class
$n_i$, $c_2(T_R|_{D_i})=c_2(S)+c_1(S)n_i$ and
$c_1(T_R|_{D_i})=c_1(S)+n_i$. The two normal classes on
the unglued normalization are $m,-m$, where $m=c_1(L)$.
Their sum is zero; $c_1(S)=F$ and $F^2=0$. The displayed
number is therefore $24-24+0=0$. Because
$\nu_*[R]=[Q]$, this is the intrinsic virtual number on $Q$.
The projective-bundle calculation is independently checked
in \texttt{ab\_push\_local\_checks.py}, using
$H(H+m)=0$, $d_0=H+m$, $d_\infty=H$.

The untwisted quotient away from its cusp has Euler number zero:
scalar multiplication supplies a free circle action, including
the two additive places before logarithmic transformation. Indeed
a circle scalar identified with an iterate of $u\phi$ must be
the identity, since every positive iterate strictly contracts
the Hermitian norm and every negative iterate expands it.
Compact fibres with a free circle action have Euler number zero.
The chosen Mumford fibre has two zero-dimensional torus orbits
and all other torus strata have Euler number zero, so
$e(X(u))=2$. The same number follows independently from the
semistable nearby-cycle count in A.4.3 below. Lemma~\ref{push:A31}
therefore excludes such a smoothing of $Q$. No requirement
on its leading smoothing section, and no extension of $f$,
entered this argument.

There is also a direct check that zero is the correct smooth
number. Replace $\phi$ in the global construction of A.1.2
by a section of $J$ vanishing on a smooth elliptic fibre of
$S$. The compact collar construction still applies. At its
zero divisor the fixed nonzero-parameter equation is $xy=uz$
with $dz\ne0$, so every nearby threefold is smooth. The
double-locus nearby Euler function is one on that elliptic
curve and zero off it; its integral is zero. This is an actual
proper-flat smoothing of $Q$, but is not Engel's family.
\end{proof}

\begin{proposition}[A.3.3; Refuted and corrected]\label{push:A33}
The assertion in Lead 3, or A.3(iii), that the two natural
fillings exhaust smooth-total-space reduced-normal-crossings
fillings is false, even when joint flatness over
$B\times\Delta$ is imposed. There are infinitely many such
fillings of the same punctured untwisted family. The
two-model classification in A.2.1 is a classification of the
specified minimal toroidal refinements, not of all fillings.
\end{proposition}
\begin{proof}
We give centres that preserve joint flatness, rather than
blowing up a point over a point of the two-dimensional base.
The section $2P$ is disjoint from $P$ and $O$. In the
Weierstrass equation its coordinates are
\[
 x(2P)=\frac{9t^2}{8},\qquad
 y(2P)=\frac{\sqrt2\,t^2(5t-32)}{32}.
\]
At $0$ it meets the nonidentity component different from
that met by $P$; at $1$ the displayed coordinates are finite
and different from those of $P,O$; at infinity the minimal
coordinates are finite and different as well. This also
checks directly the zero intersection numbers
$P\cdot2P=O\cdot2P=0$. The rational section of
$\mathcal O(P-O)$ therefore restricts to a nowhere-vanishing
holomorphic section on $2P\simeq B$. Its graph is a compact
curve $C\subset L^\times\subset R^\pm$ disjoint from both
conductor sections and from the modification centres.
It gives a curve in the smooth locus of $Q^\pm$ mapping
isomorphically to $B$.

Blow up $C$ in the smooth total fourfold $\mathcal X^\pm$.
Locally the coordinates are $(t,u,a,b)$ with
$C=(u,a,b)=0$. The blow-up charts replace $u$ by $u$ itself
or by $aU$ or $bU$. Thus the new central divisor is reduced
normal crossings, and the total space is smooth. All fibres
over $(t,u)$ still have pure dimension two; the regular base
criterion proves joint flatness. The morphism is proper and
is an isomorphism for $u\ne0$. It adds a
$\mathbb P^2$-bundle over $C$, hence a new central irreducible
component. To justify repetition, $N_{C/Q^\pm}$ is an extension
of $\mathcal O_B(-1)$ by $\mathcal O_B$ and splits, and
$H^1(N_{C/Q^\pm})=0$. The normal sequence in the total family
therefore splits off the trivial parameter direction. Its
section in the exceptional $\mathbb P^2$-bundle is disjoint
from the intersection with the strict central transform and
has the same normal sequence. We may repeat indefinitely.
Their differing component counts preclude
isomorphism as fillings. The local substitutions and the
section-doubling calculation are exact computations, recorded
with the other local checks.
\end{proof}

\subsection{A.4. Integral topology and nearby cycles}

\begin{proposition}[A.4.1; Proved]\label{push:A41}
With integral coefficients, the spaces constructed above have
\[
\begin{array}{c|rrrrrrr|r}
&b_0&b_1&b_2&b_3&b_4&b_5&b_6&\chi\\\hline
R&1&0&11&0&11&0&1&24\\
R^\pm&1&0&12&0&12&0&1&26\\
Q&1&1&8&7&10&0&1&12\\
Q^\pm&1&1&9&7&11&0&1&14
\end{array}
\]
All these homology groups are torsion free. The fundamental
groups of $R,R^\pm$ are trivial and those of $Q,Q^\pm$ are
$\mathbb Z$. The dual complex of either $Q^\pm$ is a circle.
\end{proposition}
\begin{proof}
We check the lattice step of R10 on which the calculation depends.
Let $\mathcal L=H_2(S;\mathbb Z)$, with its unimodular intersection
pairing, and put
\[
 F=[S_t],\qquad m=P-O,\qquad v=[6P]-[O]-3F.
\]
The intersections are
\[
 F^2=Fm=Fv=0,\quad m^2=-2,\quad v^2=-6,\quad mv=-1,
 \quad OF=Om=1,\quad Ov=0.
\]
Translation is the integral isometry
\[
 T_*a=a+(aF)v-(av-3aF)F.
\]
This is checked on $F,O,v$ and on the seven root classes of
the two reducible Kodaira fibres, which translation fixes.
It also gives $T^*m=m-F$. The span of $F,v,m$ is primitive:
split off the unimodular plane generated by $F,O$; the Gram
determinant of $m-F,v$ in its orthogonal complement is eleven,
so an index larger than one would have its square divide eleven.
Unimodularity therefore makes
$a\mapsto(aF,av,am)$ surjective onto $\mathbb Z^3$.

The normalization mapping cylinder gives
\[
 H_i(S)\xrightarrow{d_i}H_i(R)\longrightarrow H_i(Q)
 \longrightarrow H_{i-1}(S),\qquad
 d_i=s_{0*}-s_{\infty*}T_*.
\]
The maps are induced by actual embeddings, and the mapping
cylinder is homotopy equivalent to the pinching because the
two section embeddings are disjoint cofibrations. With $f$ the
projective-line fibre, the degree-two map is
\[
 d_2(a)=-(aF)v+(av-3aF)F+(am-aF)f.
\]
Its three coefficient columns have determinant $-1$, so the
image is primitive of rank three, the kernel has rank seven,
and the cokernel has rank eight. In degree four the image is
the primitive class Poincar\'e dual to $m$, since $mO=1$.
In degree zero the difference map is zero. The projective-bundle
formula and this exact sequence give the asserted groups of
$R,Q$. Van Kampen on the same mapping cylinder gives the
single stable loop and no other generator, hence
$\pi_1(Q)=\mathbb Z$.

Now let $h:R^+\to R$ be the modification constructed in A.2.2.
Its fibres are points off the nodal curve $C=S_\infty$, one
$\mathbb P^1$ over $C_{\mathrm{reg}}$, and a chain of two
$\mathbb P^1$'s over the node. The chart calculation identifies
the sheaf $R^2h_*\mathbb Z$ with
$\eta_*\mathbb Z_{\mathbb P^1}$, where
$\eta:\mathbb P^1\to C$ is the normalization. Here is the
integral check at the node. If $(a,b)$ are the cohomological
degrees on the two exceptional curves, generization to the
two branches is $a$ and $a+b$: on one branch the generic
exceptional fibre specializes to the first curve, on the
other it specializes to their sum. The matrix
$\left(\begin{smallmatrix}1&0\\1&1\end{smallmatrix}\right)$
is unimodular. These generizations follow by substituting
$X=za$, $v=aB$ in the blow-up charts, or by taking the two
limits of their exceptional fibres. Proper local neighbourhoods
retract onto these fibres by the positive toric contraction,
so the fibre calculation computes the higher direct images.
There is no other positive-degree direct image.

The Leray spectral sequence has only the row $H^*(R,\mathbb Z)$
and the row $H^{*-2}(\mathbb P^1,\mathbb Z)$; both are free
and in even degrees. It degenerates integrally for degree
reasons, giving the stated groups of $R^+$. The same local
neighbourhood retractions and Van Kampen show that $h$ induces
an isomorphism on fundamental groups, so $R^+$ is simply
connected. The case $R^-$ interchanges the branches.

For $Q^+$ use the normalization mapping cylinder once again.
The new degree-two difference map is the old $d_2$ with
one extra coordinate $aF$: the exceptional divisor restricts
to $F$ on the blown-up section and to zero on the other.
It has the same primitive rank-three image and the same
kernel; its cokernel rank increases by one. The degree-four
difference remains primitive, since its pushforward to $R$
already was primitive, and its cokernel likewise increases
by one. This proves the groups of $Q^+$. Interchanging the
sections changes the extra coordinate by a sign and does
not change any group. Smith normal forms are recorded in
\texttt{ab\_push\_local\_checks.out}. The dual complex has
one vertex and one self-edge, with no higher cells, and
Van Kampen again gives $\mathbb Z$.

As independent Euler checks, stratification gives
$\chi(Q)=\chi(R)-\chi(S)=24-12$. In the modification $h$,
$C_{\mathrm{reg}}\simeq\mathbb C^*$ contributes zero and
the exceptional chain above the node increases Euler number
by $3-1=2$. Thus $\chi(R^\pm)=26$ and
$\chi(Q^\pm)=26-12=14$.
\end{proof}

\begin{proposition}[A.4.2; Proved]\label{push:A42}
The untwisted smooth fibre has torsion-free integral homology
with Betti numbers
\[
 (b_0,b_1,b_2,b_3,b_4,b_5,b_6)=(1,1,8,14,8,1,1).
\]
The positive-$u$ monodromy is the identity outside degree three.
In degree three there is an integral exact sequence
\[
 0\longrightarrow \mathcal L/\langle F,v,m\rangle
 \longrightarrow H_3(X(u);\mathbb Z)
 \longrightarrow K\longrightarrow0,\qquad
 K=\{a\in\mathcal L:aF=av=am=0\},
\]
in which $T_u-I$ is zero on the left term and, on the quotient,
is the map
\[
 K\longrightarrow \mathcal L/\langle F,v,m\rangle,\qquad
 a\longmapsto a\bmod\langle F,v,m\rangle.
\]
The sign convention orients the neck circle by increasing
$\arg x$ in $xy=u$; the boundary convention for a transported
chain is chosen so that a positive $u$ loop adds this circle
times its boundary cycle. With these conventions
$(T_u-I)^2=0$, its rank is six, and its nonzero Smith factors
are $1,1,1,1,1,11$.
\end{proposition}
\begin{proof}
We use the actual semistable model A.2.2. Remove disjoint
open tubular neighbourhoods of the two conductor sections
from $R^+$, and denote the remaining compact manifold by
$Z$. Its two boundary circle bundles are identified by the
plumbing; let $E$ denote the first. The Euler class of
$E\to S$ is $n=m-F$. It is primitive. The Thom sequence
for $(R^+,Z)$ gives
\[
 H_5(Z)=\mathbb Z,\quad H_4(Z)=0,\quad
 H_3(Z)=\mathcal L/\langle F,m\rangle\simeq\mathbb Z^8,
 \quad H_2(Z)\simeq\mathbb Z^{10},\quad H_1(Z)=0.
\]
We spell out the rank and saturation assertions. If $H$ is
the projective-bundle class restricting to zero on $s_0$
and to $-m$ on $s_\infty$, and $A$ is the exceptional
divisor class, the two restriction maps on $H^2(R^+)$ are
\[
 (\alpha,h,a)\longmapsto\alpha+aF,\qquad
 (\alpha,h,a)\longmapsto\alpha-hm.
\]
Their combined image is primitive of rank twelve, and its
cokernel is $\mathcal L/\langle F,m\rangle$. In degree two,
intersection with the two sections is onto $\mathbb Z^2$:
a generic projective fibre meets both once, and an exceptional
projective fibre meets just the section containing the centre
once. In degree six the map is the diagonal
$\mathbb Z\to\mathbb Z^2$. These are precisely the Thom
maps giving the asserted groups.

The Gysin sequence of the primitive circle bundle gives
\[
 H_2(E)=n^\perp\subset\mathcal L,\qquad
 H_3(E)=\mathcal L/\mathbb Z n,\qquad
 H_1(E)=H_4(E)=0,\quad H_5(E)=\mathbb Z.
\]
The smooth fibre is obtained from $Z$ by the boundary
mapping cylinder. In its Mayer--Vietoris sequence the
degree-two difference map is the normalization difference
$d_2^+$ restricted to $n^\perp$. In the coefficient
coordinates $(s,t,h)=(aF,av,am)$ the constraint is $h=s$,
so this map has primitive rank two and kernel $K$.
The degree-three difference map is, modulo $F,m$,
$(T_*-I)a=(aF)v$, hence has primitive image $\mathbb Z v$.
These statements use the surjectivity of the three
intersection coordinates proved in A.4.1, so no integral
index is omitted. The degree-five difference is zero:
the plumbing reverses the normal circle orientation, and
the two boundary fundamental classes sum to zero in $Z$.
The degree-zero difference is zero as well. Exactness now
gives the claimed groups and the displayed degree-three
extension. Poincar\'e duality, or the same integral exact
sequence, confirms the absence of torsion in complementary
degrees.

We also compute monodromy directly from this decomposition.
In the neck $xy=u$, a positive circuit of $u$ rotates one
boundary normal coordinate once relative to the other.
A cycle in $Z$ is unchanged. A three-cycle whose cut-open
boundary represents $a\in K$ acquires the tube obtained
by sweeping $a$ once around that normal circle. The Gysin
description sends this tube to the class of $a$ in
$H_3(E)=\mathcal L/\mathbb Z n$, and then to its class in
$\mathcal L/\langle F,m,v\rangle$. This is exactly the
map in the statement. The same description in the other
degrees is zero, using the groups and difference maps just
computed. A tube has no cut-open boundary, which proves
$(T_u-I)^2=0$ without an appeal to a monodromy theorem.

Finally $K\cap\langle F,m,v\rangle=\mathbb ZF$, since
the Gram matrix of $m-F,v$ has determinant eleven.
Hence the rank is six. In the unimodular $E_8$ summand,
the primitive rank-two lattice spanned by $m-F,v$ has
discriminant eleven; its orthogonal complement maps into
the dual quotient with cokernel of order eleven. Thus the
nonzero Smith factors are five ones and eleven. The fully
integral check uses an explicit $E_8$ basis in
\texttt{ab\_push\_monodromy.py}; its output records the
kernel basis, quotient projection, and all entries of the
seven-by-seven monodromy block. This supplies a second,
exact lattice verification of the discriminant argument.
\end{proof}

\begin{proposition}[A.4.3; Proved; pointwise correction to Lead 1]\label{push:A43}
For $\mathcal X^\flat\to\Delta_u$ the nearby-fibre Euler
function on the double locus of $Q$ is zero off $S_\infty$
and one at every point of $S_\infty$, including its node.
Its integral is one. In particular the assertion that the
pointwise value on the smooth part of $S_\infty$ is zero
is false; its integrated contribution is zero.
For either semistable family $\mathcal X^\pm$ the nearby
Euler function is one off the double locus and zero on
it, giving $\chi(X(u))=26-2\cdot12=2$.
\end{proposition}
\begin{proof}
Off the zero divisor of $\phi$, the transverse nearby
fibre of $xy=u\phi$ is an annulus and has Euler number
zero. At a smooth zero, the local equation is $xy=uz$;
for fixed $u\ne0$ one solves for $z$, and the small
local nearby space is contractible. Its Euler number
is one. At the node it is the quadratic cone
$xy=uzv$, also contractible, with Euler number one.
The smooth part of the nodal elliptic curve is
$\mathbb C^*$ and has Euler number zero; the node
contributes one. The core $Q\setminus S$ is a
$\mathbb C^*$-bundle and contributes zero.
This proves the first count.

In the resolved family the local map on the double locus
is simply $ab=u$ times two transverse coordinates,
including above the old node. Its nearby annulus has
Euler number zero everywhere on that locus. Away from
it, the family is a submersion and the local value is
one. Integrating gives the second count. If the nodal
threefold in the first family is smoothed rather than
small-resolved, its ordinary-double-point Milnor fibre
is homotopy equivalent to $S^3$ and has Euler number
zero; the replacement changes $1$ to $0$, whereas
the small resolution changes $1$ to $2$. These agree
with the intrinsic virtual number in A.3.2.
\end{proof}

\begin{proposition}[A.4.4; Proved; limitation of Clemens--Schmid]\label{push:A44}
The classical K\"ahler or projective Clemens--Schmid
comparison cannot be applied to these untwisted models.
Their smooth fibres have $b_1=1$ and are not K\"ahler.
Moreover the exceptional small-resolution curve in
$R^\pm$ is null-homologous, so that normalization is
not K\"ahler either. The rational monodromy weight
filtration, defined algebraically from $N=T_u-I$ and
centred at degree three, has graded ranks $6,2,6$ in
weights $2,3,4$. This is a filtration of the computed
monodromy representation; we do not identify it with
a polarized limiting mixed Hodge structure.
\end{proposition}
\begin{proof}
The odd first Betti number excludes a K\"ahler metric.
In the generization matrix of A.4.1 a global section of
$R^2h_*\mathbb Z$ has equal branch values, hence degree
zero on the small exceptional curve. All classes pulled
back from $R$ also have degree zero on it. The Leray
calculation exhausts $H^2(R^\pm)$, so the curve has
zero homology class; it cannot have positive area for
a K\"ahler form. Finally $N^2=0$, rank $N=6$, and
$\dim H_3=14$, so the centred filtration is
$W_2=\operatorname{im}N$, $W_3=\ker N$, $W_4=H_3$.
Its graded dimensions are the asserted ones. These
topological computations replace, rather than assume,
the inapplicable K\"ahler comparison.
\end{proof}

\begin{problem}[A.4.5; Open; unrestricted scope]\label{push:A45}
Compute the integral topology and
strata of every further subdivision
in A.2.3. The calculations A.4.1--A.4.4
give the complete groups, nearby
functions and monodromy for $Q$ and
the two constructed minimal models;
they are not a topology computation
for every arbitrary refinement.
\end{problem}

\section{B. Local limits and the global component obstruction}
\subsection{B.1. The finite root neighbourhoods}

\begin{proposition}[B.1.1; Refuted and corrected]\label{push:B11}
The product assertion in Lead 4 is false on the whole root
disc. There is, however, a proper flat holomorphic filling of
the marked root family over
$\Delta_{s_j}\times\Delta_r$, where $u=r^{12}$, with smooth
total space and reduced normal-crossings $r$-central divisor.
Its central surface over an ordinary value of $s_j$ is a cycle
of twelve ruled surfaces over the smooth elliptic curve
$E'_{j,s_j}$. Their normalizations are projective bundles
associated with the degree-zero line bundle represented by
$P_{s_j}-O_{s_j}$, with the appropriate pullbacks by translations.
They are not generically products with $\mathbb P^1$.
\end{proposition}
\begin{proof}
On the root disc the elliptic family $E'_j\to\Delta_{s_j}$ is
smooth. In the additive period coordinates $(\xi,\zeta_2)$,
put $W=\exp(-2\pi i\zeta_2)$. Quotienting by
$\ker\gamma\subset\Lambda$ gives a holomorphic
$\mathbb C^*$-bundle over this elliptic family. The period
$\widehat\gamma$ acts by
\[
 (\xi,W)\longmapsto
 (\xi+6\mu(s),a(s)uW),\qquad
 a(s)=\exp(-2\pi i\beta^{\rm part}(s)).
\]
This is the bundle and translation in R8, with inverse fibre
coordinate; the period identities check its transition
functions directly. The factor $a$ is a holomorphic unit on
the entire root disc. Apply the compact-core plumbing of
A.1.2 relatively over a closed smaller root disc. Here the
smoothing section never vanishes, so the only local equation
at the conductor is $xy=u$ after a change of frame. This
gives a smooth total space before the base change.

After $u=r^{12}$, the local equation is $xy=r^{12}$,
times the two smooth coordinates on $E'_j$. Resolve its
$A_{11}$ surface singularity by its minimal toric resolution.
In lattice coordinates this inserts each primitive height-one
ray between the two boundary rays, at unit intervals. The
construction is canonical for changes $x\mapsto a x$,
$y\mapsto a^{-1}y$ by a holomorphic unit, so these local
resolutions glue. Its affine charts have the equation
$r=XY$, or $r=X$ at a smooth point of the central divisor,
times the elliptic and root coordinates. Thus the total space
is smooth, the divisor is reduced with normal crossings, and
the joint map to $(s_j,r)$ is flat. The resolution is proper,
and the plumbing was proper, so the joint map is proper.
The comparison over $r\ne0$ is the original period quotient,
since the resolution is an isomorphism there. There are
twelve components in the cycle because one application of
$\widehat\gamma$ shifts the unfolded component index by
twelve.

The ruled normalization of a component is obtained by adding
the zero and infinity sections to the indicated
$\mathbb C^*$-bundle. For a degree-zero line bundle $L_0$
on an elliptic curve, $\mathbb P(\mathcal O\oplus L_0)$
is a product over that curve only if $L_0$ is trivial.
Indeed an isomorphism would give two independent nowhere
vanishing quotient sections of a common twist of
$\mathcal O\oplus L_0$; degree zero forces that twist and
$L_0$ to be trivial. Equivalently, the two distinguished
disjoint sections have normal bundles $L_0,L_0^{-1}$.
A finite isogeny can trivialize $L_0$ only when $L_0$ is
torsion, because the norm of its pullback is a positive
tensor power of $L_0$.

The section $P$ of the rational elliptic surface is not
torsion, as is checked by its positive Mordell--Weil height
in Engel's Lemma 2.1. If its restriction to the root disc
were identically killed by a positive integer, analytic
continuation on the smooth part of $S\to B$ would give the
same identity globally. Each nonzero torsion equation
therefore has a discrete zero set. Outside a countable union
of such sets, $P_s-O_s$ is non-torsion. This disproves the
product assertion, even after an arbitrary fixed finite
isogeny of the elliptic factor. The special torsion
isogenies at $s_j=0$ used in R11 remain valid.
\end{proof}

\begin{theorem}[B.1.2; Proved]\label{push:B12}
The twisted finite deck generators extend to the models of
B.1.1. On the twelve component labels their actions are
$i\mapsto i+4$ at $j=1$ and $i\mapsto i-3$ at $j=2$,
with labels in $\mathbb Z/12$. Every nonidentity group
element acts freely, including at every node of the central
cycle. The quotient is a smooth complex fourfold, proper
and flat over $(t,r)=(s_j^{m_j},r)$, with a reduced
normal-crossings $r$-central divisor. Its fibre divisor
over $t=p_j$ has multiplicity exactly $m_j$.

The reduction over $s_j=r=0$ is the central surface of R11:
a cycle of four ruled surfaces at $j=1$, and a cycle of
three at $j=2$. In the root presentation the normal
character to $s_j=0$ is $e^{-2\pi i/m_j}$ and the
character on $r$ is trivial.
\end{theorem}
\begin{proof}
We retain source and target arguments. Write primes for
evaluation at $g_js$, put $\zeta_j=e^{-2\pi i/m_j}$,
and set
\[
 R_1=\begin{pmatrix}-1/\tau&0\\(1-\mu)/\tau&1\end{pmatrix},
 \qquad
 R_2=\begin{pmatrix}1/\tau&0\\-\mu/\tau&1\end{pmatrix}.
\]
The identity $R_j\Pi_c=\Pi'_c A_j$ and the fixed-vector
identities $A_jv_j=v_j$ are checked symbolically in
\texttt{ab\_push\_finite\_decks.py}. In the coordinates
$(s,\xi,W,r)$ the two affine deck maps become
\begin{align*}
 s'&=\zeta_1s,\qquad
 \xi'=-\xi/\tau+(6\mu'+2\tau'-4)/3,\\
 W'&=r^4 U_1(s,\xi)W,\\
 U_1&=\exp\!\left[-2\pi i
  \left(\frac{1-\mu}{\tau}\xi+
        \frac{\beta^{{\rm part}\prime}+2\mu'}3\right)\right],
 \end{align*}
and
\begin{align*}
 s'&=\zeta_2s,\qquad
 \xi'=\xi/\tau+(-6\mu'-3\tau'+3)/4,\\
 W'&=r^{-3} U_2(s,\xi)W,\\
 U_2&=\exp\!\left[-2\pi i
  \left(-\frac{\mu}{\tau}\xi-
        \frac{\beta^{{\rm part}\prime}+3\mu'}4\right)\right].
 \end{align*}
Each $U_j$ is a nowhere vanishing holomorphic function on
the additive elliptic cover. Period covariance says
exactly that it defines the required bundle isomorphism
after quotient by $\ker\gamma$; it is not being treated
as a scalar function on the elliptic curve when it is not
one. Multiplication by a unit preserves every toric chart.
Multiplication by $r^a$ sends the chart of component index
$i$ to the chart of index $i+a$. Therefore the displayed
maps and their inverses extend holomorphically to the
resolved infinite chain and then to its twelve-component
quotient. The identities $g_1^3=g_2^4=1$ on the quotient
hold on the dense punctured family, and hence everywhere.

A nonidentity power has component shift $4,8$ in the first
case and $-3,-6,-9$ in the second. None is zero modulo
twelve. The same shifts act on the edge labels of the
cycle, so no node can be fixed either. At $s\ne0$ the
rotation of $s$ already excludes fixed points. The finite
quotient map is consequently a holomorphic local
isomorphism. Its source is smooth, so its target is smooth.
Properness follows from the proper root model and the
finite base map. Flatness can be checked after the
surjective local-isomorphism cover: there $s^{m_j}$ and
$r$ are a regular sequence, and their fibres have pure
dimension two. The local central equation remains
$r=XY$ or $r=X$. Finally $t=s^{m_j}$ shows that the
Cartier divisor over $p_j$ is exactly $m_j$ times its
reduction, including where this divisor meets a central
node. This also proves the asserted normal characters.

At $s=0$ use the explicitly marked isogenies
\[
 L_1(\xi,\eta)=(\eta,\xi-\eta/3),\qquad
 L_2(\xi,\eta)=(\eta,\xi-\eta/2).
\]
Their kernel generators are $(2/3,\rho)$ and $(1/2,i)$,
and the finite deck maps lift to
\[
 (\xi,\eta)\mapsto(\xi+\kappa_1/3,(\rho-1)\eta),
 \qquad
 (\xi,\eta)\mapsto(\xi-\kappa_2/4,-i\eta).
\]
In the Tate coordinate $v=e^{-2\pi i\xi}$ these are the
shifts $4,-3$ just computed. The kernel generators act
by nonzero translations on the smooth elliptic factor
$F_j$, hence freely even at the two fixed points of each
Tate component. Every other element with nonzero deck
power has a nonzero component shift. The full product
groups are therefore free on the product semistable
models. Their component orbits have cardinalities three
and four, giving precisely four and three ruled
components. The original lattice isogenies intertwine
the maps on the dense tori and their unit-interval fans,
so the quotient identifications extend across all nodes.
This verifies the particular comparison with R11 that
we use.

For the punctured family obtained by descending the
$c$-family, the holomorphic correction $\eta_j(s)$ in
R2 causes no additional local obstruction here. At the
shift $c\mapsto c+12$, the reference translations
$-4\Pi_c(\widehat u+\widehat w)$ and
$3\Pi_c(\widehat u+\widehat w)$ are periods.
The remaining correction is equivariant and independent
of $c$. Translation by $-c\eta_j(s)/12$ on the additive
cover conjugates that descent to the period descent.
It commutes with the finite deck action by equivariance.
Thus the constructed finite models fill the actual local
pullback by $w=r^6$ after a specified holomorphic
identification over $r\ne0$. This local gauge is not
asserted to extend across $r=0$ on overlaps with other
base neighbourhoods.
\end{proof}

\subsection{B.2. The base-changed cusp fan}

\begin{proposition}[B.2.1; Proved]\label{push:B21}
For the plain base change $u=r^n$ of the untwisted cusp
model, the smooth period-invariant refinements which
retain the prescribed bottom triangulation, are jointly
flat over $(q,r)$, and have reduced $r$-central divisor
are indexed by $k=0,\ldots,n$. In a fundamental prism
write its four bottom rays as $A,B,C,D$ as in A.2.1,
and put $\rho_i=(0,i,0,1)$ for $0\leq i\leq n$.
The maximal cones for the index $k$ are
\[
 ABC\rho_k,\quad BCD\rho_k,\quad
 AB\rho_i\rho_{i+1}\ (i<k),\quad
 CD\rho_i\rho_{i+1}\ (i\geq k).
\]
There are thirteen such models for $n=12$ and three
for $n=2$. Each has smooth total space and a reduced
normal-crossings central divisor with $n$ components
and no triple locus. Its component labels over the
punctured $q$-disc are globally $\mathbb Z/n$; in
particular their monodromy around $q=0$ is the identity.
\end{proposition}
\begin{proof}
The map of cocharacter lattices for the base change is
\[
 (y_1,y_2,h_q,h_r)\longmapsto
 (y_1,y_2,h_q,nh_r).
\]
The old top ray indexed by $l$ pulls back to
$(0,nl,0,1)$. Smoothness on the two-dimensional top
face forces insertion of every primitive ray
$\rho_{nl+i}$, $0<i<n$: a skipped interval of length
$d>1$ would retain lattice index $d$ on that face.
Joint flatness excludes any new ray with both base
heights positive, and the fixed bottom excludes new
rays of $h_r=0$. Reducedness permits only top rays
of $h_r=1$. These are consequently all the rays.

On the left face $AC\rho_0\rho_n$, a triangulation
with these collinear top vertices has a unique triangle
containing the bottom edge $AC$. Write its third
vertex as $\rho_a$. The triangles on either side
of that bridge are then forced. The analogous right
face has bridge $BD\rho_b$. The first period takes
$A,C$ to $B,D$ and fixes every top ray, so periodicity
requires $a=b=k$. The fixed bottom triangles $ABC$
and $BCD$ must then cone to $\rho_k$. The remaining
successive top intervals force exactly the displayed
cones, leaving no choice. Conversely those cones
have determinant $\pm1$ and common-face
intersections. Their normalized volumes sum to
$n+2$, the volume of the prism, so they cover it.
The determinant, intersection and volume assertions
for both requested exponents are certified in
\texttt{ab\_push\_basechange\_fan.py} using rational
arithmetic.

The period actions are
\[
 G_1:(y_1,y_2)\mapsto(y_1+h_q,y_2),\qquad
 G_2:(y_1,y_2)\mapsto(y_1,y_2+h_q+nh_r).
\]
Thus the choices propagate consistently from the
fundamental prism. The proper toric subdivisions
descend by the covering construction of A.1.2.
Every ray has pure $q$-height or pure $r$-height, so
in a regular affine chart $q$ and $r$ are monomials
in disjoint sets of coordinates. This proves joint
flatness. All nonzero $r$-heights equal one, and
there are at most two top rays in any maximal cone;
this proves reduced normal crossings and absence of
triple points. Properness follows by descent from
the proper toric modification of the base-changed
proper cusp model.

The central components and all their intersection
maps are specified by the stars of the $\rho_i$.
Modulo the periods their indices are $i\bmod n$;
adjacent indices meet in the conductor elliptic
surface. Away from the cusp their normalizations
are the ruled bundles of the expanded plumbing.
At the unique bridge class the star has both bottom
triangles, and gives the nodal-curve blow-up and
small resolution described in A.2.2, on the
corresponding ruled bundle. This star description
also fixes the gluing when the bridge is on an
inserted component. The endpoint choices $k=0,n$
are the canonical resolutions of the base changes
of the two original models. Intermediate values
are related to them by the successive side flops.

Finally, $G_1$ fixes each index and $G_2$ changes it
by $n$. The units in the cusp periods are
single-valued holomorphic functions of $q$ and do
not change a ray label. The toric charts themselves
are defined over $q$, not a root of $q$. Hence
these labels remain global on the punctured
$q$-disc, proving the monodromy assertion. Further
smooth periodic subdivisions with mixed rays give
infinitely many models flat over $r$ alone, as in
A.2.1; they are excluded by the stated joint-flatness
condition.
\end{proof}

\begin{problem}[B.2.2; Open; unrestricted scope]\label{push:B22}
Classify and describe all further
period-invariant smooth subdivisions
of the base-changed cusp fans when
joint flatness is not required.
B.2.1 gives all jointly flat reduced
models, but the unrestricted class
again contains infinitely many further
refinements.
\end{problem}

\subsection{B.3. The prescribed pieces do not glue}

\begin{theorem}[B.3.1; Refuted; Lead 5]\label{push:B31}
The finite models of B.1 and any of the plain
$u=r^{12}$ cusp models of B.2 cannot be glued to a
proper complex space over $B\times\Delta_r$ with
the prescribed central fibres on all three base
neighbourhoods. In particular Lead 5's asserted
simultaneous gluing is false. This conclusion does
not rule out a different cusp model of the same
punctured $S^6$ family.
\end{theorem}
\begin{proof}
Suppose such a gluing exists and put
$B^\circ=B\setminus\{0,1,\infty\}$. Normalize
its reduced central fibre over $B^\circ$. In
each of the prescribed plumbing charts this is
a proper smooth family whose fibres are the
disjoint union of twelve connected ruled
surfaces. The charts and their transition
isomorphisms therefore give a locally constant
twelve-element set of connected components.
Equivalently, the Stein factor of this
normalization is an unramified degree-twelve
cover of $B^\circ$: the algebra of fibrewise
holomorphic functions is locally $\mathcal O^{12}$,
since each proper connected ruled surface has
only constant holomorphic functions. This
description constructs the cover and requires
no conjecture about components at the missing
three points.

The double curves give the adjacency graph of
a twelve-cycle. Thus its transition
permutations lie in the dihedral group
$D_{12}$, the automorphism group of that graph.
We record a peripheral permutation by its
deck-identification map. Holonomy defined by
lifting a clockwise path and returning its
endpoint to the initial root slice uses the
inverse permutations throughout; the relation
and the following contradiction are unchanged.
At $0$ the monodromy is a rotation by four
steps; at $1$ it is a rotation by minus three
steps, by the explicit root deck maps in
B.1.2. Relabelling a cycle can reverse either
orientation, but changes these only to
$\pm4$ and $\pm3$. At the cusp the
monodromy is the identity by B.2.1.

The peripheral relation in
$\pi_1(B^\circ)$ is $g_1g_2g_0=1$ with
the consistently chosen clockwise loops.
In any common labelling, the product of the
two finite rotations would therefore have
to be the identity. Its possible rotation
numbers are
\[
 4+3,\quad4-3,\quad-4+3,\quad-4-3
 \equiv7,1,11,5\pmod {12}.
\]
None is zero. This contradicts the peripheral
relation. The four residues and the finite
component cycles are checked exactly by
\texttt{ab\_push\_finite\_decks.py}.
\end{proof}

\begin{proposition}[B.3.2; Proved; necessary condition only]\label{push:B32}
If a corrected global construction retains the
finite models and a twelve-cycle over
$B^\circ$, then, in the common orientation
with finite shifts $4,-3$, its cusp
component monodromy must be the twelve-cycle
of shift $-1$. The associated component
cover of $B^\circ$ is connected. Its compact
normalization over $B$ has degree twelve,
ramification indices $3,4,12$, and genus
three. In this deck convention it is the
normalization of
\[
 y^{12}=\frac{t^4}{(t-1)^3}.
\]
These assertions do not construct
the required central complex space.
\end{proposition}
\begin{proof}
The peripheral relation gives the remaining
rotation as $-4+3=-1$. The subgroup generated
by $4,-3$ is all of $\mathbb Z/12$, proving
connectedness. The local cycle decompositions
have four cycles of length three, three
cycles of length four, and one of length
twelve. Filling a punctured-disc cover by
its normal branched cover gives ramification
contribution
$4(3-1)+3(4-1)+(12-1)=28$.
The Riemann--Hurwitz formula gives
$2-2g=24-28=-4$, hence $g=3$.
The displayed Kummer equation has orders
$4,-3,-1$ at the three points and hence
the same marked cyclic monodromy. The
classification of connected covering
spaces identifies its restriction to
$B^\circ$ with the required component
cover; normalization gives the compact
comparison. Reversing the transport
convention replaces $y$ by $y^{-1}$.
Only this component cover has been
constructed; no topology is assigned here
to a hypothetical threefold above it.
\end{proof}

\begin{proposition}[B.3.3; Proved; obstruction to a specified repair]\label{push:B33}
A cyclic cusp-root construction using a
fractional translation by the two cusp
periods cannot repair B.3.1 while retaining
the reduced Mumford fibre for $r\ne0$.
More precisely, let $q=s^d$, $u=r^n$,
and suppose a cyclic deck generator has
translation $d^{-1}(aG_1+bG_2)$ in the
additive notation for the two periods.
Assume its $r$-component shift is integral.
Then its $r$-component shift is
$nb/d\bmod n$, whereas the multiplicity
of every generic quotient $q$-component
is a positive multiple of
\[
 \frac{d}{\gcd(d,a,b)}.
\]
A nonzero $r$-component shift therefore
forces a nonreduced $q$-fibre. In
particular a shift $-1\bmod12$ in this
class of constructions forces multiplicity
at least twelve.
\end{proposition}
\begin{proof}
In the coordinates of A.1.2 the two
periods have $q$-valuation vectors
$(1,0)$ and $(0,1)$; their $r$-valuation
vectors are $(0,0)$ and $(0,n)$.
After $q=s^d$, the period lattice on
the $q$-central component indices is
$d\mathbb Z^2$. A regular toroidal
resolution has indices in
$(\mathbb Z/d)^2$. The proposed deck
generator translates them by $(a,b)$.
The order of that translation is
$d/\gcd(d,a,b)$, so the subgroup
preserving a component has order
$h=\gcd(d,a,b)$. Its subgroup fixing
a generic point has some order
$h_0$ dividing $h$; unit factors may
make $h_0$ smaller. That isotropy
acts on the normal coordinate $s$
by a primitive $h_0$th root of
unity. The quotient normal
coordinate is $s^{h_0}$, and
$q=(s^{h_0})^{d/h_0}$ has
multiplicity $d/h_0$, a positive
multiple of $d/h$. Unit factors
cannot decrease this lower bound.

Reducedness would require $h_0=d$,
and therefore $h=d$,
so both $a,b$ are divisible by $d$.
Then $nb/d$ is zero modulo $n$.
If $nb/d$ is relatively prime to
$n$, its order in $\mathbb Z/n$
is $n$, which divides the order
$d/h$ of the full component
translation; this gives the final
lower bound. The residue calculations
and examples for the requested
exponents are independently recorded
in \texttt{ab\_push\_cusp\_root\_repair.py}.
This rules out the specified repair,
not a cusp model with a different
geometry or a different comparison
to the punctured family.
\end{proof}

\begin{theorem}[B.3.4; Proved; split-normalization trace obstruction]\label{push:B34}
Let $C$ be the genus-three component cover
in B.3.2, and let $\mathscr E\to C$ be
the relatively minimal elliptic surface
obtained by pulling back $S\to B$.
Write $\sigma$ for the order-twelve deck
transformation, with its natural lift to
$\mathscr E$. Then
\[
 \operatorname{Tr}_{\sigma}
   \mathscr E(\mathbb C(C))
 :=\left\{\sum_{i=0}^{11}\sigma^i Q:
                 Q\in\mathscr E(\mathbb C(C))\right\}
 =12\mathbb Z P.
\]
In particular there is no section $Q$
for which $(t_Q\circ\sigma)^{12}=t_{6P}$.
Thus the proposed repair that takes the
normalization over $C$ to be ruled over
this split elliptic surface, with adjacent
conductor map $t_Q\circ\sigma$, is
impossible.

This conclusion is restricted to the
stated split normalization. A ruled
normalization over a nontrivial genus-one
torsor need not admit a section or such
an expression for its affine deck lift;
it is not excluded by this theorem.
\end{theorem}
\begin{proof}
There is a convenient model of the cover
and its elliptic surface:
\[
 C:\ b^3=a^4+1,\qquad
 t=b^3,\quad y=b/a,\qquad
 \mathscr E:\ Y^2=X^3-3bX+2a^2.
\]
Indeed $y^{12}=t^4/(t-1)^3$.
The substitution from the original
Weierstrass coordinates is
$x=a^2b^4X$, $y_{\rm old}=a^3b^6Y$.
The affine discriminant of the new
equation is $1728$; it has no singular
fibre on the affine curve. The curve
$C$ has one point at infinity, with
pole orders $3,4$ for $a,b$.
Put $z=a/b$ and $c=z^3a=z^4b$ there.
Then $z$ is a uniformizer and
\[
 c^3(1-c)=z^{12},\qquad c(0)=1.
\]
In the integral Weierstrass coordinates
$X_0=z^2X$, $Y_0=z^3Y$, the equation is
\[
 Y_0^2=X_0^3-3cX_0+2c^2.
\]
Its discriminant has order twelve and
its $c_4$ is a unit. The unique singular
fibre of the minimal regular surface
is therefore $I_{12}$. The fundamental
line bundle is $\mathcal O_C(\infty)$,
so $\chi(\mathcal O_{\mathscr E})=1$.
This is a projective elliptic surface.

Let $\omega=e^{2\pi i/3}$ and
$\lambda=e^{2\pi i/12}$. The deck lift is
\[
 \sigma(a,b,X,Y)=(ia,\omega b,\lambda^2X,\lambda^3Y).
\]
It fixes the section $P=(0,-\sqrt2a)$.
Two further sections are
\[
 R=\left(\frac{a^2}{b-1},
                 \frac{\sqrt3a}{b-1}\right),
 \qquad
 R'=\sigma R=
   \left(\frac{a^2}{b-\omega},
                 \frac{\sqrt3\omega a}{b-\omega}\right).
\]
Direct substitution verifies their
equations. The three orbit points
$R,\sigma R,\sigma^2R$ lie on the line
\[
 Y=\frac{\sqrt3b}{a}X-\sqrt3a.
\]
Thus $R+\sigma R+\sigma^2R=O$ by the
elliptic chord law, and $\sigma^3R=R$.
The ansatz that produced these sections,
its other solutions, and this line
identity are checked symbolically in
\texttt{ab\_push\_trace\_sections.py}.

We give the intersection computation
needed for the height lattice; this
also fixes the component indices used
below. The only pole of $R$ on the affine
curve is at $(a,b)=(0,1)$. There $a$ is
a uniformizer and $b-1=a^4/3+O(a^8)$,
so $X=3a^{-2}+O(a^2)$ and
$Y=3\sqrt3a^{-3}+O(a)$.
Consequently $R\cdot O=1$.
The same computation at $(0,\omega)$
gives $R'\cdot O=1$, while $P\cdot O=0$.
At the other two points with $a=0$,
$P$ and $R$ meet simply: the fibre is
smooth at $(X,Y)=(0,0)$, and the
difference of their $Y$ coordinates
has a nonzero linear term in $a$.
Thus $P\cdot R=2$, and likewise
$P\cdot R'=2$. The sections $R,R'$
meet simply at $(a,b)=(0,\omega^2)$.
They cannot meet elsewhere on the
affine curve, since their $X$
coordinates are distinct for $a\ne0$.

To check infinity, for the section
$\sigma^jR$ put $k=\omega^j$. Its
integral coordinates are
\[
 X_0=\frac{c^2}{c-kz^4},\qquad
 Y_0=\frac{\sqrt3kc z^4}{c-kz^4}.
\]
At the node use analytic coordinates
\[
 U=Y_0-(X_0-1)\sqrt{X_0+2},\qquad
 V=Y_0+(X_0-1)\sqrt{X_0+2}.
\]
They satisfy
\[
 UV=(1-c)\{3X_0-2(c+1)\},
\]
whose right side is a unit times $z^{12}$.
The section expansions are
\[
 U=-\frac{\sqrt3k^2}{6}z^8+O(z^{12}),
 \qquad V=2\sqrt3kz^4+O(z^8).
\]
Hence $R,R'$ both meet component $4$
of the resolved $I_{12}$ cycle, with
distinct coordinates on that component.
Reversing the cycle numbering replaces
both indices by $8$ and has no effect
on the argument. The section $P$ meets
the identity component, at the smooth
point $(X_0,Y_0)=(0,-\sqrt2)$.
Thus no further intersections occur at
infinity, and $R\cdot R'=1$.

The height pairing formula and its
positivity are those of
\cite[Theorems 8.4 and 8.6, table (8.16)]{Shioda90}.
For an $I_{12}$ fibre its correction is
$\operatorname{contr}(i,j)=\min(i,j)-ij/12$.
The preceding intersections give, in
the ordered basis $P,R,R'$,
\[
 G=\begin{pmatrix}
 2&0&0\\
 0&4/3&-2/3\\
 0&-2/3&4/3
 \end{pmatrix},\qquad \det G=8/3.
\]
In particular these are three independent
sections. The Leray sequence and
$R^1\pi_*\mathcal O_{\mathscr E}
=\mathcal O_C(-\infty)$ show that its
irregularity equals $g(C)=3$; the
right-hand sheaf has no global section.
This also agrees with the Picard-variety
statement in \cite[Theorem 4.1]{Shioda90}.
The equalities $\chi=1$ and
$e(\mathscr E)=12$ give $p_g=3$,
$b_2=22$ and $h^{1,1}=16$.
The zero section, a fibre and the eleven
nonidentity $I_{12}$ components generate
a rank-thirteen trivial lattice.
The identification of sections with
the N\'eron--Severi quotient
\cite[Theorem 1.3]{Shioda90}, together
with $\rho\le h^{1,1}$, now shows that
the Mordell--Weil rank is exactly three
and $\rho=16$.

There is no nonzero torsion section.
Indeed its height would be zero, so
for some $0\le j<12$ the height formula
would give
\[
 0=2+2(Q\cdot O)-\frac{j(12-j)}{12}.
\]
The correction is at most three and
$Q\cdot O$ is a nonnegative integer.
This forces $Q\cdot O=0$ and
$j(12-j)=24$, which has no integral
solution.

Let $L=\mathbb ZP+\mathbb ZR+\mathbb ZR'$.
The integral N\'eron--Severi sublattice
generated by the trivial lattice and
these three sections has determinant
$-32$. This follows either by eliminating
the trivial lattice with the displayed
height matrix, or from the integral
intersection matrix recorded in
\texttt{ab\_push\_trace\_obstruction.py}.
All its entries have just been computed:
sections have square $-1$, their mutual
intersections are $2,2,1$, and $R,R'$
meet precisely component $4$.
If $d=[\mathscr E(\mathbb C(C)):L]$,
the section--N\'eron--Severi quotient
identification gives the same index
for this sublattice in
$\operatorname{NS}(\mathscr E)$.
Its integral determinant implies
$d^2\mid32$, hence $d\in\{1,2,4\}$.
In particular $4Q\in L$ for every section.

In the rational basis $P,R,R'$ the
deck action and its trace are
\[
 M=\begin{pmatrix}1&0&0\\0&0&-1\\0&1&-1\end{pmatrix},
 \qquad \sum_{i=0}^{11}M^i=
       \operatorname{diag}(12,0,0).
\]
Suppose first that the trace of $Q$ is
$6P$. By the index bound we may write
\[
 Q=\frac12P+\frac14(uR+vR'),\qquad u,v\in\mathbb Z.
\]
Let $j\in\mathbb Z/12$ be its component
index at infinity. Since both $R,R'$
meet component four, the height formula
modulo integers gives
\[
 \frac{2u-v}{6}\equiv\frac j3\pmod{\mathbb Z},
 \qquad
 \frac{-u+2v}{6}\equiv\frac j3\pmod{\mathbb Z}.
\]
Equivalently $2u-v-2j$ and $-u+2v-2j$
are divisible by six. Reducing modulo
two shows that both $u,v$ are even.
Write $u=2u'$, $v=2v'$. The component
map is a group homomorphism, so
\[
 2j\equiv4(u'+v')\pmod{12}.
\]
Therefore $j$ is even. The geometric
height formula then lies in
$\tfrac13\mathbb Z$. In contrast the
height from $G$ is
\[
 \langle Q,Q\rangle
 =\frac12+\frac{u'^2-u'v'+v'^2}{3},
\]
which lies in the disjoint coset
$\tfrac12+\tfrac13\mathbb Z$.
This is a contradiction.

The same argument determines the full
trace image. If $Q=\alpha P+\beta R+\gamma R'$
in the rational span, pairing with the
narrow section $P$ gives
$2\alpha=\langle Q,P\rangle\in\mathbb Z$.
Every odd half-integral value of $\alpha$
leads to the identical congruences and
height contradiction, with $1/2$
replaced by an odd square divided by
two. Thus $\alpha\in\mathbb Z$ and
the trace is in $12\mathbb ZP$.
The section $P$ realizes the generator.
The exact matrices, local series,
torsion test, and all relevant residue
classes are verified in
\texttt{ab\_push\_trace\_obstruction.py}.
Finally the identity
$(t_Q\circ\sigma)^{12}=t_{\operatorname{Tr}_\sigma Q}$
proves the claimed obstruction, with
the scope stated in the theorem.
\end{proof}

\begin{corollary}[B.3.5; Proved; exact Mordell--Weil lattice]\label{push:B35}
In the notation of B.3.4,
\[
 \mathscr E(\mathbb C(C))
   =\mathbb ZP\oplus\mathbb ZR\oplus\mathbb ZR',
 \qquad \det\operatorname{NS}(\mathscr E)=-32.
\]
\end{corollary}
\begin{proof}
The affine coordinate ring of $C$ is
$\mathbb C[a,b]/(b^3-a^4-1)$, with
unique representatives of degree at most
two in $b$. The pole orders $3i+4j$ of
$a^ib^j$, $0\le j\le2$, are distinct,
so cancellation cannot create functions
with pole order one or two at infinity.
Consequently
\[
 H^0(C,\mathcal O_C(2\infty))=\mathbb C,
 \qquad
 H^0(C,\mathcal O_C(3\infty))
       =\mathbb C\oplus\mathbb Ca.
\]
A section disjoint from $O$ therefore
has $X=x_0$, $Y=y_0+y_1a$ for constants
$x_0,y_0,y_1$. Substitution in the
Weierstrass equation forces $x_0=y_0=0$
and $y_1^2=2$. Thus the only nonzero
sections with intersection zero with
$O$ are $P,-P$.
Every other nonzero section has
intersection at least one with $O$,
and hence height at least $4-3=1$.

B.3.4 proves that every section has
integral coefficient along $P$ in the
rational height decomposition. The
remaining lattice can have index only
one, two or four over $\mathbb ZR+\mathbb ZR'$.
Any proper such overlattice contains a
nonzero element of order two modulo
this sublattice. After subtracting an
integral combination of $R,R'$ it is
one of $R/2$, $R'/2$, $(R+R')/2$.
Each has height $1/3$, contradicting
the preceding lower bound. Thus the
index is one. The determinant already
computed in B.3.4 is the full
N\'eron--Severi determinant.
The pole-basis calculation, coefficient
equations and three half-class heights
are included in the exact certificate
\texttt{ab\_push\_trace\_obstruction.py}.
\end{proof}

\begin{proposition}[B.3.6; Proved; an unramified torsor alternative]\label{push:B36}
Put $K=\mathbb C(t)$, $L=\mathbb C(C)$,
$G=\operatorname{Gal}(L/K)=\langle\sigma\rangle$,
and let $E/K$ be the generic elliptic
curve of $S\to B$. There is a nonzero
class
\[
 \xi\in\operatorname{Sha}(E_L/L)^G
 \quad\hbox{whose transgression in}\quad
 H^2(G,E(L))=\mathbb ZP/12\mathbb ZP
 \quad\hbox{is }6P.
\]
Here $\operatorname{Sha}$ denotes the
subgroup of genus-one torsors trivial
over the henselian local fields of all
points of $C$.
In particular the trace obstruction of
B.3.4 does not obstruct all genus-one
torsors. Such a class has a relatively
minimal projective genus-one surface
$V\to C$, locally isomorphic to
$\mathscr E\to C$, and an affine lift
$A:V\to V$ of $\sigma$ with
\[
 A^{12}=t_{6P}.
\]
This is an existence statement for a
possible normalization base, not a
construction of a threefold central
fibre or of its smoothing.
\end{proposition}
\begin{proof}
We use Galois cohomology with its
continuous cochains. The five-term
Hochschild--Serre sequence gives
\[
 H^1(K,E)\longrightarrow H^1(L,E)^G
 \xrightarrow{\ d_2\ }H^2(G,E(L))
 \longrightarrow H^2(K,E).
\]
The transgression sequence is
\cite[Chapter III, Section 4, Theorem 2
and the following remark, pp.129--130]{HS53};
its continuous form follows by passing
to finite Galois quotients for the
finitely many values of each continuous
cochain. For this complex function field,
$H^2(K,E)=0$, as in
\cite[Lemma 4.7.8]{CDL24}.
One can also see this from
$\operatorname{cd}K=1$: the torsion
submodule of $E(\overline K)$ has no
second cohomology, and its quotient is
a rational vector space, whose positive
continuous cohomology vanishes by
averaging on finite quotients.

Cyclic group cohomology gives
$H^2(G,E(L))=E(K)/\operatorname{Tr}_\sigma E(L)$.
By B.3.4 this is $\mathbb Z/12$, with
generator the class of $P$. Hence there
is a $G$-invariant torsor class
$\xi_0\in H^1(L,E)$ with $d_2\xi_0=6P$.
We next remove all its local ramification
without changing this transgression.

At a place $v$ of $B$ choose a point
$w$ of $C$ above it and its decomposition
subgroup $G_w$. The only nontrivial
decomposition groups have orders
$3,4,12$, at $0,1,\infty$ respectively.
The component groups of the original
fibres are $\mathbb Z/3$, $\mathbb Z/2$
and zero. Thus $6P$ belongs to the
identity component of the local N\'eron
group at every ramified place.
Multiplication by the ramification
index is surjective on that identity
group over the henselian local field:
on the residue group it is surjective
on $\mathbb G_a(\mathbb C)$ or
$\mathbb G_m(\mathbb C)$, and its
derivative is invertible; Hensel's
lemma lifts a division point.
Choose such a division point over
$K_v^h$. Viewed over $L_w^h$, it is
invariant, so its local norm is $6P$.
The restriction of the transgression
class to $H^2(G_w,E(L_w^h))$ is
therefore zero. For unramified places
the decomposition group is trivial,
so the same assertion is automatic.

Naturality and exactness of the local
five-term sequences now show that
$(\xi_0)_w$ is the restriction of some
$\eta_v\in H^1(K_v^h,E)$.
Only finitely many of these classes
need be nonzero. The original elliptic
surface is nontrivial, since
$j=1728t$. Ogg--Shafarevich localization
is therefore surjective:
\[
 H^1(K,E)\longrightarrow
 \bigoplus_{v\in B}H^1(K_v^h,E)
 \longrightarrow0.
\]
This is the combination of exact
sequence (4.7.4) and Theorem 4.7.9
of \cite{CDL24}; the same surjectivity
and its geometric fibre-twist
interpretation are given in
\cite[Section 4, equation (4)]{Xiao91}.
Take a global $\eta$ with these local
values and put
$\xi=\xi_0-\operatorname{res}_{L/K}\eta$.
It is locally trivial at every point
of $C$, remains $G$-invariant, and has
$d_2\xi=6P$. It is nonzero because its
transgression is nonzero.

The standard interpretation of the
Weil--Ch\^atelet group represents
$\xi$ by a smooth projective genus-one
curve over $L$; taking its relatively
minimal regular model over $C$ gives
$V$. Local triviality identifies each
henselian model with that of $E_L$.
Over $\mathbb C$ the corresponding
\'etale identifications yield
holomorphic identifications over small
base discs. Equivalently, $V$ is obtained
by gluing copies of $\mathscr E$ over
these discs by translations by local
sections. These translations extend
over the regular $I_{12}$ model.
Thus $V$ is locally isomorphic to
$\mathscr E$, has smooth total surface,
no multiple fibres, and the same single
$I_{12}$ fibre. Properness and flatness
hold locally over the base and hence
globally. This use of torsors and local
twists is the one in
\cite[Sections 3--4]{Xiao91}; it does
not assert an explicit equation for
the selected torsor.

The $G$-invariance of $\xi$ gives an
affine isomorphism of its generic fibre
covering $\sigma$. It extends to the
relatively minimal regular surface:
in this particular situation the
extension can be checked on a common
resolution. If an exceptional
$(-1)$-curve for one projection is not
contracted by the other, its image is
a vertical rational curve and has
self-intersection at least $-1$,
since blowing down only increases
self-intersection. On either locally
trivial model all vertical rational
curves are components of the $I_{12}$
fibre, of square $-2$. Thus every such
$(-1)$-curve is contracted by both
projections and can be removed from
the common resolution. Repeating this
proves that the map is an isomorphism.
The same argument gives the local
comparison with the Jacobian model
used above. Transporting the base by
$\sigma$ changes none of these facts.
Write $A_0$ for this lift. Its twelfth
power is translation by a section
$N\in E(L)^G=\mathbb ZP$.
Choosing the powers $1,A_0,\ldots,A_0^{11}$
as lifts of $G$ gives the cyclic
factor set
$N\lfloor(i+j)/12\rfloor$.
Its cohomology class is precisely the
descent obstruction $d_2\xi$, so
$N=(6+12m)P$ for an integer $m$.
Replacing $A_0$ by $A=t_{-mP}\circ A_0$
gives $A^{12}=t_{6P}$.
The finite carry-cocycle identities,
its exact order two modulo twelve,
and the three local component checks
are recorded in
\texttt{ab\_push\_torsor\_cocycles.py}.

Neither this cohomological construction
nor local triviality specifies the
normal line bundles, the threefold
cusp charts, or a smoothing with the
given punctured-family identification.
Those remain requirements of B.4.1.
\end{proof}

\subsection{B.4--B.5. The remaining filling problem}

\begin{problem}[B.4.1; Open]\label{push:B41}
Construct, or obstruct, a proper flat
holomorphic filling of the actual $S^6$
family after a finite base change which
repairs the cusp incompatibility of B.3.1.
If it exists, determine an equivariant
model for descent to $w$, its central
multiplicities and a resolution, and the
least semistable base-change exponent.
The local constructions in B.1 and B.2
do not settle this problem.
\end{problem}

The obstruction cannot be removed by citing
$H^1(B,\mathcal V)=0$ from R2. That
vanishing solves an additive translation
cocycle for a fixed nonzero parameter and
its integer shifts. On a finite overlap
the translation contains
\[
 \frac{\log s_j}{2\pi i}\,\Pi_c v_j,
 \qquad c=-\frac{12}{2\pi i}\log r.
\]
Its vertical exponential contains a product
of the two logarithms. No assertion in
the cited vanishing controls its growth
at $r=0$ or makes it a unit in the central
toric charts. B.3.1 proves that at least
one of those extension requirements must
fail for the prescribed pieces. A
construction with a corrected cusp remains
necessary before descent or characteristic
numbers can be computed.

\begin{problem}[B.5.1; Open]\label{push:B51}
For an actual corrected filling in B.4.1,
compute the integral topology of its
central fibre and of a descent resolution,
the nearby-cycle Euler count, and the
applicable monodromy or Hodge comparison.
These invariants are not assigned to the
incompatible assembly refuted in B.3.1.
\end{problem}
